%% 第8巻　空間に広がる量　本文（第14回編集委員会）
%% 章は \bkinputchap{ch8-NN} で読む（まだない章は、組むときに飛ばして「BKPAGE missing」を出す）
%% 第8巻は巻の色が濃い紺なので、部扉の絵の線は白（bkOn）でかく
%% Ver2（全巻にかかる修正事項 No.58、組版のメモリ対策）：第9巻と同じく、章の前に Lua のごみ集め（collectgarbage）をして、使っているメモリをログに書く
\providecommand\bkinputchap[1]{\directlua{collectgarbage("collect") texio.write_nl("log","BKMEM #1 "..math.floor(collectgarbage("count")/1024).." MB")}\IfFileExists{#1.tex}{\input{#1}}{\clearpage\typeout{BKPAGE missing #1}}}
\bkinputchap{ch8-00}

%% ======================== 第1部 たくさんの変数の微分 ========================
\bkPartOpener{1}{第8巻\hspace{0.8em}第1部}{たくさんの変数の微分}%
  {172/{山の地図から、坂の急さを読みとれる？},
   160/{東へ進むときと北へ進むときで、坂の急さはちがう？},
   148/{ボールは、どの向きに転がり出す？}}%
  {% 等高線と、いちばん急な向きの矢印（白い線）
   \begin{scope}[bkOn, line width=1.1pt, line cap=round, line join=round]
     \draw (144,124) ellipse (6 and 4);
     \draw (144,124) ellipse (13 and 8.5);
     \draw (144,124) ellipse (21 and 13.5);
     \draw[-{Stealth[length=2.8mm]}, line width=1.6pt] (165,124) -- (172,124);
     \draw[-{Stealth[length=2.8mm]}, line width=1.6pt] (144,137.5) -- (144,143);
   \end{scope}}%
  {場所を決めると高さや温度が1つ決まる関数を、地図の等高線で見るよ。1つの向きだけに動いたときの変わり方（偏微分）を集めると、いちばん急な向きが矢印でわかるよ。}
\bkPartContents{1}{%% 自動生成（gen.py）
\bktocchap{8-1}{地図と等高線：2つ以上の変数の関数}{場所を決めると、高さが1つ決まる}{1}{4}{4}{30}{C}
\bktocchap{8-2}{偏微分：ほかを止めて1つだけ動かす}{東へ進むときの坂と、北へ進むときの坂}{1}{4}{4}{64}{C}
\bktocchap{8-3}{全微分：少し動いたときの変化}{ななめに少し進むと、高さはどれだけ変わる？}{1}{5}{5}{100}{C}
\bktocchap{8-4}{いちばん急な向き：勾配 grad}{ボールが転がる向きは、等高線と直角}{1}{5}{5}{138}{C}
}%
{2つ以上の変数の関数を等高線と表で読み、偏微分・全微分・勾配を計算できるようになります。ロボットのヤコビ行列（第7巻）の作り方もここでわかります。}
\bkpartbegin{1}{たくさんの変数の微分}
\bkinputchap{ch8-01}
\bkinputchap{ch8-02}
\bkinputchap{ch8-03}
\bkinputchap{ch8-04}

%% ======================== 第2部 ベクトル場と積分 ========================
\bkPartOpener{2}{第8巻\hspace{0.8em}第2部}{ベクトル場と積分}%
  {172/{空間のどこにも矢印がある、ってどういうこと？},
   160/{遠回りすると、した仕事は変わる？},
   148/{あみを通りぬける水の量は、どう数える？}}%
  {% 矢印の場と、それを通りぬける面（白い線）
   \begin{scope}[bkOn, line width=1.1pt, line cap=round, line join=round]
     \foreach \y in {110,120,130,140} {\foreach \x in {118,130,142} {\draw[-{Stealth[length=2.2mm]}] (\x,\y) -- ++(8,{1.5*sin(\y*6)});}}
     \draw[line width=1.6pt] (156,108) -- (166,114) -- (166,144) -- (156,138) -- cycle;
   \end{scope}}%
  {空間の点ごとに矢印が1本ずつある「ベクトル場」を、道にそって・面を通して、たし合わせるよ。積分の名前と、実際にふつうの定積分に直す手順を、場合ごとにていねいに見ていくよ。}
\bkPartContents{2}{%% 自動生成（gen.py）
\bktocchap{8-5}{ベクトル場}{空間の点ごとに、矢印を1本ずつ}{6}{4}{4}{176}{C}
\bktocchap{8-6}{面や立体でたし合わせる：重積分}{厚さがちがう板の重さを求める}{2}{5}{5}{202}{C}
\bktocchap{8-7}{道にそってたす：線積分}{遠回りすると、仕事は変わる？}{2}{5}{5}{252}{C}
\bktocchap{8-8}{面を通りぬける量：面積分と流束}{あみの向きで、通りぬける水の量が変わる}{2}{5}{5}{292}{C}
}%
{ベクトル場を図と式で表し、重積分・線積分・面積分を、場合ごとの手順でふつうの定積分に直して計算できるようになります。}
\bkpartbegin{2}{ベクトル場と積分}
\bkinputchap{ch8-05}
\bkinputchap{ch8-06}
\bkinputchap{ch8-07}
\bkinputchap{ch8-08}

%% ======================== 第3部 発散と回転 ========================
\bkPartOpener{3}{第8巻\hspace{0.8em}第3部}{発散と回転}%
  {172/{蛇口から広がる水の流れを、数で表すには？},
   160/{風車が回るかどうかは、何で決まる？},
   148/{中のわき出しと、外へ出ていく量の関係は？}}%
  {% わき出しと、うず（白い線）
   \begin{scope}[bkOn, line width=1.1pt]
     \fill (136,126) circle (1.8);
     \foreach \a in {0,45,...,315} \draw[-{Stealth[length=4pt]}] (136,126) ++(\a:4) -- ++(\a:8);
     \draw[dash pattern=on 2.5pt off 1.5pt] (130,120) rectangle (142,132);
     \foreach \r in {7,12} {\draw[-{Stealth[length=4pt]}] (164,122) ++(30:\r) arc (30:150:\r);
                            \draw[-{Stealth[length=4pt]}] (164,122) ++(210:\r) arc (210:330:\r);}
     \foreach \a in {0,90,180,270} \fill (164,122) -- ++(\a:3.2) -- ++(\a+50:1.8) -- cycle;
     \fill (164,122) circle (0.7);
   \end{scope}}%
  {この部では、流れの中に小さな箱や風車を置いて、場所ごとの「わき出し」と「うず」を数で表すよ。電気・磁気や水の流れの式を読むための、大事な道具だよ。}
\bkPartContents{3}{%% 自動生成（gen.py）
\bktocchap{8-9}{わき出しと吸いこみ：発散 div}{小さな箱から、どれだけ出ていく？}{6}{5}{5}{350}{C}
\bktocchap{8-10}{うず：回転 rot}{小さな風車を置くと、回るかな？}{6}{5}{5}{396}{C}
\bktocchap{8-11}{中と外をつなぐ：ガウスの定理とストークスの定理}{中のわき出しを全部たすと、外へ出る量になる}{2}{5}{5}{440}{C}
\bktocchap{8-12}{ならす働き：ラプラシアン}{まわりの平均との差が、変化を決める}{1}{5}{5}{492}{C}
}%
{発散・回転・ラプラシアンの意味を、小さな箱と風車で説明し、ガウスの定理・ストークスの定理で「中」と「外」をつなげられるようになります。}
\bkpartbegin{3}{発散と回転}
\bkinputchap{ch8-09}
\bkinputchap{ch8-10}
\bkinputchap{ch8-11}
\bkinputchap{ch8-12}

%% ======================== 第4部 偏微分方程式 ========================
\bkPartOpener{4}{第8巻\hspace{0.8em}第4部}{偏微分方程式}%
  {172/{金属の棒に、熱はどう広がる？},
   160/{波の形は、どうやって伝わる？},
   148/{電気と磁気の4つの式は、何を言っている？}}%
  {% 熱の広がり（時間とともに平らになる山、白い線）
   \begin{scope}[bkOn, line width=1.1pt, line cap=round]
     \draw[line width=0.6pt] (116,110) -- (172,110);
     \draw[line width=1.6pt, domain=-26:26, samples=60] plot ({144+\x}, {110+28*exp(-\x*\x/40)});
     \draw[line width=1.1pt, domain=-26:26, samples=60] plot ({144+\x}, {110+16*exp(-\x*\x/120)});
     \draw[line width=0.8pt, dash pattern=on 2pt off 2pt, domain=-26:26, samples=60] plot ({144+\x}, {110+10*exp(-\x*\x/300)});
   \end{scope}}%
  {場所と時間の両方で変わる量のルールが、偏微分方程式だよ。熱の広がりや波の伝わり方を表計算で再現し、電気と磁気の4つの式を「わき出し」と「うず」で読むよ。}
\bkPartContents{4}{%% 自動生成（gen.py）
\bktocchap{8-13}{熱と波の伝わり方：偏微分方程式}{金属の棒に熱が広がるようすを、表計算で再現する}{1}{5}{5}{528}{C}
\bktocchap{8-14}{div と rot で書かれた式を読もう}{電気と磁気の4つの式を、「わき出し」と「うず」で読む}{6}{5}{5}{568}{R}
}%
{熱方程式・波動方程式を表計算で解き、マクスウェル方程式を div と rot のことばで読めるようになります。第9巻と物理シリーズの電磁気・流体の土台です。}
\bkpartbegin{4}{偏微分方程式}
\bkinputchap{ch8-13}
\bkinputchap{ch8-14}

%% ======================== 巻末 ========================
\bkinputchap{ch8-back}
